\textbf{Phần 1. Trắc nghiệm lựa chọn}
\setcounterpageref{ex}{0}

\begin{ex} %[2D10N49]
Số nào sau đây là số tròn trăm?
\choice
{$150$}
{$205$}
{\True $200$}
{$99$}
\loigiai{
Số tròn trăm tận cùng là hai chữ số 0: 200.\\
\textbf{Đáp án của câu hỏi là C}
}
\end{ex}

\begin{ex} %[2D10N49]
Số nào sau đây là số tròn chục?
\choice
{\True $320$}
{$325$}
{$305$}
{$321$}
\loigiai{
Số tròn chục tận cùng là chữ số 0: 320.\\
\textbf{Đáp án của câu hỏi là A}
}
\end{ex}

\begin{ex} %[2D10H49]
Số tròn trăm liền sau số 400 là
\choice
{$401$}
{\True $500$}
{$300$}
{$410$}
\loigiai{
Số tròn trăm liền sau của 400 là 500.\\
\textbf{Đáp án của câu hỏi là B}
}
\end{ex}

\begin{ex} %[2D10H49]
Số tròn chục liền trước số 150 là
\choice
{$149$}
{$160$}
{\True $140$}
{$100$}
\loigiai{
Số tròn chục liền trước là $150 - 10 = 140$.\\
\textbf{Đáp án của câu hỏi là C}
}
\end{ex}

\begin{ex} %[2D10H49]
Các số tròn trăm từ 100 đến 500 là
\choice
{$100, 200, 400, 500$}
{$100, 150, 200, 250$}
{$200, 300, 400, 500$}
{\True $100, 200, 300, 400, 500$}
\loigiai{
Gồm 5 số: 100, 200, 300, 400, 500.\\
\textbf{Đáp án của câu hỏi là D}
}
\end{ex}

\begin{ex} %[2D10H49]
Tính: $300 + 200 = \dots$
\choice
{\True $500$}
{$400$}
{$600$}
{$50$}
\loigiai{
$300 + 200 = 500$.\\
\textbf{Đáp án của câu hỏi là A}
}
\end{ex}

\begin{ex} %[2D10V49]
Tính: $800 - 300 = \dots$
\choice
{$400$}
{$600$}
{\True $500$}
{$50$}
\loigiai{
$800 - 300 = 500$.\\
\textbf{Đáp án của câu hỏi là C}
}
\end{ex}

\begin{ex} %[2D10V49]
Tính: $200 + 400 + 100 = \dots$
\choice
{\True $700$}
{$600$}
{$800$}
{$900$}
\loigiai{
$200 + 400 + 100 = 700$.\\
\textbf{Đáp án của câu hỏi là A}
}
\end{ex}

\begin{ex} %[2D10V49]
Số tròn chục lớn nhất có ba chữ số là
\choice
{$900$}
{\True $990$}
{$999$}
{$1000$}
\loigiai{
Số tròn chục lớn nhất có 3 chữ số là 990.\\
\textbf{Đáp án của câu hỏi là B}
}
\end{ex}

\begin{ex} %[2D10V49]
Số tròn trăm lớn nhất có ba chữ số là
\choice
{$990$}
{$999$}
{$1000$}
{\True $900$}
\loigiai{
Số tròn trăm lớn nhất có 3 chữ số là 900.\\
\textbf{Đáp án của câu hỏi là D}
}
\end{ex}

\begin{ex} %[2D10V49]
Có bao nhiêu số tròn trăm có ba chữ số?
\choice
{$8$ số}
{\True $9$ số (từ $100$ đến $900$)}
{$10$ số}
{$11$ số}
\loigiai{
Từ 100 đến 900 có 9 số tròn trăm.\\
\textbf{Đáp án của câu hỏi là B}
}
\end{ex}

\begin{ex} %[2D10C49]
Tìm một số tròn trăm biết rằng lấy số đó cộng với 300 rồi trừ đi 200 thì được 600.
\choice
{$400$}
{$600$}
{$700$}
{\True $500$}
\loigiai{
$600 + 200 - 300 = 500$.\\
\textbf{Đáp án của câu hỏi là D}
}
\end{ex}

\textbf{Phần 2. Trắc nghiệm đúng sai}
\setcounterpageref{ex}{0}

\begin{ex} %[2D10H49]
Xét tính đúng sai về số tròn trăm:
\choiceTF
{\True Số tròn trăm có chữ số hàng chục và hàng đơn vị đều là 0.}
{\True $100$ là số tròn trăm bé nhất có ba chữ số.}
{$990$ là số tròn trăm lớn nhất có ba chữ số.}
{\True $900$ là số tròn trăm lớn nhất có ba chữ số.}
\loigiai{
\begin{itemize}
    \item Tận cùng là 00: Đúng.
    \item 100 là số tròn trăm bé nhất: Đúng.
    \item 990 là số tròn chục (không phải tròn trăm). Ý 3 Sai, Ý 4 Đúng.
\end{itemize}
\begin{itemize}
    \item a) \textbf{Đúng}.
    \item b) \textbf{Đúng}.
    \item c) \textbf{Sai}.
    \item d) \textbf{Đúng}.
\end{itemize}
}
\end{ex}

\begin{ex} %[2D10V49]
Cho các phép tính với số tròn trăm:
\choiceTF
{\True $400 + 300 = 700$.}
{\True $900 - 500 = 400$.}
{$200 + 600 = 900$.}
{\True $200 + 600 = 800$.}
\loigiai{
\begin{itemize}
    \item $400 + 300 = 700$: Đúng.
    \item $900 - 500 = 400$: Đúng.
    \item $200 + 600 = 800 \neq 900$. Ý 3 Sai, Ý 4 Đúng.
\end{itemize}
\begin{itemize}
    \item a) \textbf{Đúng}.
    \item b) \textbf{Đúng}.
    \item c) \textbf{Sai}.
    \item d) \textbf{Đúng}.
\end{itemize}
}
\end{ex}

\textbf{Phần 3. Trả lời ngắn}
\setcounterpageref{ex}{0}

\begin{ex} %[2D10H49]
Tính: $400 + 500$.
\shortans{900}
\loigiai{
\textbf{Đáp số:} $900$\\
$400 + 500 = 900$.
}
\end{ex}

\begin{ex} %[2D10H49]
Tính: $700 - 300$.
\shortans{400}
\loigiai{
\textbf{Đáp số:} $400$\\
$700 - 300 = 400$.
}
\end{ex}

\begin{ex} %[2D10V49]
Tìm số tròn trăm liền trước số $600$.
\shortans{500}
\loigiai{
\textbf{Đáp số:} $500$\\
$600 - 100 = 500$.
}
\end{ex}

\begin{ex} %[2D10V49]
Tìm số tròn trăm liền sau số $700$.
\shortans{800}
\loigiai{
\textbf{Đáp số:} $800$\\
$700 + 100 = 800$.
}
\end{ex}

\begin{ex} %[2D10V49]
Tính giá trị biểu thức: $1000 - 400 - 200$.
\shortans{400}
\loigiai{
\textbf{Đáp số:} $400$\\
$600 - 200 = 400$.
}
\end{ex}

\begin{ex} %[2D10C49]
Có bao nhiêu số tròn trăm nằm giữa số 200 và số 700? (Gồm các số: 300, 400, 500, 600)
\shortans{4}
\loigiai{
\textbf{Đáp số:} $4$\\
Có 4 số.
}
\end{ex}
